\noindent\textbf{Por qué se usa este procedimiento.} La ecuación característica del bloque plano da media $c=(\sigma_x+\sigma_y)/2$ y radio $R=\sqrt{[(\sigma_x-\sigma_y)/2]^2+\tau_{xy}^2}$. Las dos principales son $c+R$ y $c-R$; se conserva además la principal fuera del plano cuando corresponda.

\begin{enumerate}

\item El radio vale $50$ MPa, un resultado exacto.

\end{enumerate}

\noindent\textbf{Operaciones y resultados intermedios.}

\begin{enumerate}

\item Tensión principal mayor.

\[\fitmath{y_{1}=\frac{60+0}{2}+\sqrt{{\left(\frac{60-0}{2}\right)}^{2}+{40}^{2}}\approx 80\,\mathrm{MPa}}\]

\end{enumerate}
